% ==============================================================================
% Appendix: Formal Proofs
"Holonomy Geometry of Cognitive States in Transformer Internal Representations"
% ==============================================================================

\section*{Appendix A: Formal Proofs}
\label{app:proofs}

This appendix provides complete, self-contained proofs for the key propositions and
theorems stated in Section~\ref{sec:instantiation}. Each proof is presented in full
detail to facilitate verification and to establish the logical foundations on which the
main results rest.

\paragraph{Self-citation note.}
The companion papers Zhang (2025a, 2025b) formalize the abstract geometric and logical
foundations of the LCESA framework and are currently under preparation for concurrent
submission. All definitions and propositions cited from these works are reproduced in full
in the present paper for self-contained readability; no result in the present paper
depends on the companion papers for its proof.

We recall the standing notation: $d$ is the residual stream dimension,
$K = \{\mathrm{min}, \mathrm{nar}, \mathrm{soc}, \mathrm{mor}, \mathrm{pos}\}$ is the set
of standpoint layers, $d_k = \dim(W_k)$ for each $k \in K$ with $\sum_k d_k = d$, and
$P_{W_k}$ denotes the orthogonal projector onto the subspace $W_k$. All transport
operators $U_{ij}$ are assumed to satisfy the block-diagonal structure of
Proposition~\ref{prop:block-structure}.

% ==============================================================================
\subsection{Proof of Proposition: Metric Preservation of Attention Transport (Orthogonal Special Case)}
\label{app:proof-metric-preservation}

\begin{proposition}[Metric Preservation of Attention Transport, Orthogonal Case]
\label{prop:metric-preservation-app}
Under the orthogonal subspace assumption ($W_k \perp W_l$ for $k \neq l$, the $\delta = 0$
special case of Assumption~\ref{ass:independence}) and $\mathbf{V}_h^\top \mathbf{V}_h =
\mathbf{I}_{d_v}$ for all $h$, the transport operator $\mathbf{U}_{ij} = \sum_h
\alpha_{ji}^{(h)} \mathbf{V}_h \mathbf{V}_h^\top$ satisfies:
\[
    \mathbf{U}_{ij}^\top \mathbf{g}_F \mathbf{U}_{ij} = \mathbf{g}_F
\]
when $\alpha_{ji}^{(k)} = 1$ for all $k$ (normalized attention).
\end{proposition}

\begin{proof}
We proceed in three steps: first we decompose $U_{ij}$ into block form, then we
evaluate the quadratic form $U_{ij}^\top g_F U_{ij}$ block-by-block, and finally
we specialize to the case where all attention weights are unity.

\medskip\noindent\textbf{Step 1: Block decomposition of $U_{ij}$.}
By Proposition~\ref{prop:block-structure}, the transport operator admits the
decomposition
\[
    U_{ij} = \sum_{k \in K} \bar{\alpha}_{ji}^{(k)}\, P_{W_k},
\]
where $\bar{\alpha}_{ji}^{(k)} = \frac{1}{|H_k|} \sum_{h \in H_k} \alpha_{ji}^{(h)}$ is
the mean attention weight of layer $k$'s heads from $x_i$ to $x_j$, and $P_{W_k}$ is the
orthogonal projector onto $W_k$. This decomposition follows from the orthonormality of the
value matrices: for each head $h \in H_k$, the matrix $V_h V_h^\top$ is the orthogonal
projector onto $\operatorname{range}(V_h)$ (since $V_h^\top V_h = I_{d_v}$ implies $V_h
V_h^\top$ is idempotent and symmetric), and $\sum_{h \in H_k} V_h V_h^\top = P_{W_k}$
because the ranges of $\{V_h\}_{h \in H_k}$ form an orthonormal basis for $W_k$.

\medskip\noindent\textbf{Step 2: Evaluation of $U_{ij}^\top g_F U_{ij}$.}
Since $U_{ij}$ is a linear combination of orthogonal projectors onto mutually orthogonal
subspaces, it is symmetric:
\[
    U_{ij}^\top = \Big(\sum_{k \in K} \bar{\alpha}_{ji}^{(k)}\, P_{W_k}\Big)^\top
               = \sum_{k \in K} \bar{\alpha}_{ji}^{(k)}\, P_{W_k}^\top
               = \sum_{k \in K} \bar{\alpha}_{ji}^{(k)}\, P_{W_k} = U_{ij}.
\]
The fiber metric (Definition~\ref{def:fiber-metric}) is
\[
    g_F = \bigoplus_{k \in K} \alpha_k\, I_{d_k} = \sum_{k \in K} \alpha_k\, P_{W_k},
\]
where the $\alpha_k > 0$ are the layer weights. Now compute:
\begin{align*}
    U_{ij}^\top\, g_F\, U_{ij}
    &= U_{ij}\, g_F\, U_{ij}
      && \text{(since } U_{ij} = U_{ij}^\top \text{)} \\[4pt]
    &= \Big(\sum_{k \in K} \bar{\alpha}_{ji}^{(k)}\, P_{W_k}\Big)\,
       \Big(\sum_{l \in K} \alpha_l\, P_{W_l}\Big)\,
       \Big(\sum_{m \in K} \bar{\alpha}_{ji}^{(m)}\, P_{W_m}\Big) \\[4pt]
    &= \sum_{k,l,m \in K} \bar{\alpha}_{ji}^{(k)}\, \alpha_l\,
       \bar{\alpha}_{ji}^{(m)}\, P_{W_k}\, P_{W_l}\, P_{W_m}.
\end{align*}
By the orthogonality of the subspaces (the $\delta = 0$ special case of
Assumption~\ref{ass:independence}),
$P_{W_k} P_{W_l} = \delta_{kl}\, P_{W_k}$. Therefore:
\[
    P_{W_k}\, P_{W_l}\, P_{W_m} = \delta_{kl}\, \delta_{lm}\, P_{W_k} = \delta_{klm}\,
    P_{W_k},
\]
where $\delta_{klm} = 1$ if $k = l = m$ and $0$ otherwise. The triple sum collapses:
\begin{equation}
\label{eq:metric-quadratic}
    U_{ij}^\top\, g_F\, U_{ij}
    = \sum_{k \in K} \big(\bar{\alpha}_{ji}^{(k)}\big)^2\, \alpha_k\, P_{W_k}.
\end{equation}

\medskip\noindent\textbf{Step 3: Specialization to $\alpha_{ji}^{(k)} = 1$ for all $k$.}
When $\alpha_{ji}^{(h)} = 1$ for every head $h$, we have $\bar{\alpha}_{ji}^{(k)} = 1$ for
all $k \in K$. Moreover, when $\alpha_k = 1$ for all $k$ (the standard Euclidean fiber
metric, as stated in the convention following Definition~\ref{def:fiber-metric}),
equation~\eqref{eq:metric-quadratic} becomes:
\[
    U_{ij}^\top\, g_F\, U_{ij}
    = \sum_{k \in K} 1^2 \cdot 1 \cdot P_{W_k}
    = \sum_{k \in K} P_{W_k}
    = I_d.
\]
The last equality holds because $\{W_k\}_{k \in K}$ is an orthogonal decomposition of
$\mathbb{R}^d$ (Proposition~\ref{prop:fiber-decomposition}), so $\sum_k P_{W_k} = I_d$.
Since $g_F = \sum_k P_{W_k} = I_d$ under the standard metric, we conclude
\[
    U_{ij}^\top\, g_F\, U_{ij} = I_d = g_F,
\]
which establishes that $U_{ij}$ is an isometry of the fiber metric under normalized
attention.

\medskip\noindent\textbf{Remark.} In the general case where $\alpha_k$ are not all equal
to $1$, equation~\eqref{eq:metric-quadratic} shows that $U_{ij}^\top g_F U_{ij} =
\sum_k (\bar{\alpha}_{ji}^{(k)})^2 \alpha_k P_{W_k}$, which equals $g_F = \sum_k
\alpha_k P_{W_k}$ if and only if $(\bar{\alpha}_{ji}^{(k)})^2 = 1$ for all $k$. Since
attention weights are non-negative, this requires $\bar{\alpha}_{ji}^{(k)} = 1$ for all
$k$, which is the normalized attention condition stated in the proposition.
\end{proof}

% ==============================================================================
\subsection{Proof of Theorem: Curvature as Path-Dependence}
\label{app:proof-path-dependence}

For completeness, we reproduce the proof of Theorem~\ref{thm:path-dependence} with
additional detail on the algebraic manipulations and the geometric interpretation of each
direction.

\begin{theorem}[Curvature as Path-Dependence, restated from Theorem~\ref{thm:path-dependence}]
\label{thm:path-dependence-app}
$F_{ijk} = I_d$ if and only if the transport from $x_i$ to $x_k$ is the same whether or
not the intermediate event $x_j$ intervenes:
\[
    F_{ijk} = I_d \quad\Longleftrightarrow\quad U_{ij} \cdot U_{jk} = U_{ik}^{\exp}.
\]
\end{theorem}

\begin{proof}
The holonomy deviation (Definition~\ref{def:curvature}) is defined by
\[
    F_{ijk} = U_{ij} \cdot U_{jk} \cdot \big(U_{ik}^{\exp}\big)^{-1}.
\]
We note that $U_{ik}^{\exp} \in \operatorname{GL}(d, \mathbb{R})$ is invertible by
definition: it is the expected transport operator, estimated from baseline conversations
in which standpoint coherence is preserved. The block-diagonal structure
(Proposition~\ref{prop:block-structure}) ensures that $U_{ik}^{\exp} = \bigoplus_k
\bar{\alpha}^{(k)}_{\exp}\, I_{d_k}$ with each $\bar{\alpha}^{(k)}_{\exp} > 0$, so
invertibility is guaranteed.

\medskip\noindent\textbf{Forward direction ($\Rightarrow$):} Suppose $F_{ijk} = I_d$.
By Definition~\ref{def:curvature}:
\[
    U_{ij} \cdot U_{jk} \cdot \big(U_{ik}^{\exp}\big)^{-1} = I_d.
\]
Since $U_{ik}^{\exp}$ is invertible, we may right-multiply both sides by
$U_{ik}^{\exp}$:
\[
    U_{ij} \cdot U_{jk} \cdot \big(U_{ik}^{\exp}\big)^{-1} \cdot U_{ik}^{\exp}
    = I_d \cdot U_{ik}^{\exp},
\]
which simplifies to
\[
    U_{ij} \cdot U_{jk} = U_{ik}^{\exp}.
\]
This establishes that the composite transport through the intermediate event $x_j$ equals
the expected direct transport from $x_i$ to $x_k$. Geometrically, the connection transport
along the path $x_i \to x_j \to x_k$ yields the same result as the expected flat
transport along $x_i \to x_k$: the intermediate event $x_j$ did not alter the standpoint,
and the transport is path-independent.

\medskip\noindent\textbf{Reverse direction ($\Leftarrow$):} Suppose $U_{ij} \cdot U_{jk} =
U_{ik}^{\exp}$. Substituting into the definition of curvature:
\begin{align*}
    F_{ijk}
    &= U_{ij} \cdot U_{jk} \cdot \big(U_{ik}^{\exp}\big)^{-1} \\
    &= U_{ik}^{\exp} \cdot \big(U_{ik}^{\exp}\big)^{-1}
      && \text{(by hypothesis)} \\
    &= I_d.
\end{align*}
Hence the curvature is trivial: the standpoint was transported consistently, with no
path-dependence introduced by the event $x_j$.

\medskip\noindent\textbf{Algebraic summary.} The two directions establish the logical
equivalence
\[
    F_{ijk} = I_d
    \quad\Longleftrightarrow\quad
    U_{ij} \cdot U_{jk} = U_{ik}^{\exp},
\]
which we may restate as: the holonomy around the triangle $\langle x_i, x_j, x_k
\rangle$ is trivial if and only if the composite transport through $x_j$ coincides with
the expected flat transport.
\end{proof}

% ==============================================================================
\subsection{Proof of Proposition: Gauge Invariance of Block-Specific Curvature}
\label{app:proof-gauge-invariance}

We provide a self-contained proof of Proposition~\ref{prop:gauge-invariance} with full
detail on the use of Frobenius norm invariance under orthogonal conjugation.

\begin{proposition}[Gauge Invariance of Block-Specific Curvature, restated from
Proposition~\ref{prop:gauge-invariance}]
\label{prop:gauge-invariance-app}
Under the gauge transformation $U_{ij} \mapsto g_j\, U_{ij}\, g_i^{-1}$ with $g_i, g_j
\in G = \prod_{k \in K} O(d_k)$, the block-specific curvature norms are invariant:
\[
    \big\|g_i\, [F_{ijk}]_k\, g_i^{-1} - \Pi_k\big\|_F
    = \big\|[F_{ijk}]_k - \Pi_k\big\|_F
    \qquad \forall\, k \in K,
\]
where $\Pi_k$ is the oblique projector onto $W_k$ (Definition~\ref{def:oblique-projector})
and $[F_{ijk}]_k = \Pi_k\, F_{ijk}\, \Pi_k$.
\end{proposition}

\begin{proof}
We first establish the block structure of the gauge-transformed curvature, then prove
the norm invariance using the unitary invariance of the Frobenius norm.

\medskip\noindent\textbf{Step 1: Block structure of the gauge-transformed curvature.}
Under a gauge transformation, each transport operator transforms as $U_{ij} \mapsto g_j\,
U_{ij}\, g_i^{-1}$. The composite transport transforms as:
\begin{align*}
    U_{ij} \cdot U_{jk}
    &\;\longmapsto\; (g_j\, U_{ij}\, g_i^{-1}) \cdot (g_k\, U_{jk}\, g_j^{-1})
    = g_k\, (U_{ij} \cdot U_{jk})\, g_i^{-1}.
\end{align*}
The expected flat transport transforms in the same way:
\[
    U_{ik}^{\exp} \;\longmapsto\; g_k\, U_{ik}^{\exp}\, g_i^{-1}.
\]
Therefore the curvature transforms by conjugation:
\begin{align*}
    F_{ijk}
    &\;\longmapsto\; g_k\, (U_{ij} \cdot U_{jk})\, g_i^{-1}
       \cdot \big(g_k\, U_{ik}^{\exp}\, g_i^{-1}\big)^{-1} \\
    &= g_k\, (U_{ij} \cdot U_{jk})\, g_i^{-1}
       \cdot g_i\, (U_{ik}^{\exp})^{-1}\, g_k^{-1} \\
    &= g_k\, (U_{ij} \cdot U_{jk})\, (U_{ik}^{\exp})^{-1}\, g_k^{-1} \\
    &= g_k\, F_{ijk}\, g_k^{-1}.
\end{align*}

\medskip\noindent\textbf{Step 2: Block-level conjugation.}
Since $g_k \in G = \prod_{l \in K} O(d_l)$, we write $g_k = (g_k^{(\mathrm{min})}, \ldots,
g_k^{(\mathrm{pos})})$ with each $g_k^{(l)} \in O(d_l)$. The element $g_k$ acts on
$\mathbb{R}^d$ as the block-diagonal matrix
\[
    g_k = \bigoplus_{l \in K} g_k^{(l)} = \operatorname{diag}\!\big(g_k^{(\mathrm{min})},\,
    g_k^{(\mathrm{nar})},\, \ldots,\, g_k^{(\mathrm{pos})}\big).
\]
Under the block-diagonal structure of $F_{ijk}$ (which follows from the block-diagonal
structure of each $U_{ij}$ by Proposition~\ref{prop:block-structure}, since the product
and inverse of block-diagonal matrices are block-diagonal), the conjugation $g_k\,
F_{ijk}\, g_k^{-1}$ acts blockwise:
\[
    [g_k\, F_{ijk}\, g_k^{-1}]_l = g_k^{(l)}\, [F_{ijk}]_l\, \big(g_k^{(l)}\big)^{-1}
    \qquad \forall\, l \in K.
\]

\medskip\noindent\textbf{Step 3: Frobenius norm invariance under orthogonal conjugation.}
We now prove the key algebraic identity: for any $Q \in O(n)$ (an $n \times n$ orthogonal
matrix) and any $A \in \operatorname{Mat}(n, \mathbb{R})$,
\begin{equation}
\label{eq:frob-invariance}
    \|Q A Q^\top\|_F = \|A\|_F.
\end{equation}
To establish this, recall that the Frobenius norm satisfies $\|M\|_F^2 =
\operatorname{tr}(M^\top M)$ for any matrix $M$. Therefore:
\begin{align*}
    \|Q A Q^\top\|_F^2
    &= \operatorname{tr}\!\big((Q A Q^\top)^\top\, (Q A Q^\top)\big) \\
    &= \operatorname{tr}\!\big(Q A^\top Q^\top\, Q A Q^\top\big)
      && \text{(transpose of a product)} \\
    &= \operatorname{tr}\!\big(Q A^\top\, (Q^\top Q)\, A Q^\top\big) \\
    &= \operatorname{tr}\!\big(Q A^\top A Q^\top\big)
      && \text{(since } Q^\top Q = I_n \text{)} \\
    &= \operatorname{tr}\!\big(A^\top A\, Q^\top Q\big)
      && \text{(cyclic property of trace: } \operatorname{tr}(XY) = \operatorname{tr}(YX)\text{)} \\
    &= \operatorname{tr}\!\big(A^\top A\big)
      && \text{(since } Q^\top Q = I_n \text{)} \\
    &= \|A\|_F^2.
\end{align*}
Taking square roots (both sides are non-negative) yields identity~\eqref{eq:frob-invariance}.

\medskip\noindent\textbf{Step 4: Application to the block-specific curvature norm.}
For each layer $l \in K$, the gauge-transformed block-specific curvature is
\[
    [F_{ijk}^{\,\mathrm{gauge}}]_l = g_k^{(l)}\, [F_{ijk}]_l\, \big(g_k^{(l)}\big)^{-1}.
\]
Since $g_k^{(l)} \in O(d_l)$, we have $(g_k^{(l)})^{-1} = (g_k^{(l)})^\top$, so this is
an orthogonal conjugation. The block-specific curvature norm is:
\begin{align*}
    \big\|[F_{ijk}^{\,\mathrm{gauge}}]_l - I_{d_l}\big\|_F
    &= \big\|g_k^{(l)}\, [F_{ijk}]_l\, (g_k^{(l)})^{-1} - I_{d_l}\big\|_F.
\end{align*}
Now observe that $g_k^{(l)}\, \Pi_l\, (g_k^{(l)})^{-1} = \Pi_l$, since $g_k^{(l)} \in
O(d_l)$ maps $W_l$ to itself and $\Pi_l$ depends only on the subspace $W_l$ (not on the
choice of basis). Therefore:
\begin{align*}
    g_k^{(l)}\, [F_{ijk}]_l\, (g_k^{(l)})^{-1} - \Pi_l
    &= g_k^{(l)}\, [F_{ijk}]_l\, (g_k^{(l)})^{-1}
       - g_k^{(l)}\, \Pi_l\, (g_k^{(l)})^{-1} \\
    &= g_k^{(l)}\, \big([F_{ijk}]_l - \Pi_l\big)\, (g_k^{(l)})^{-1}.
\end{align*}
Applying identity~\eqref{eq:frob-invariance} with $Q = g_k^{(l)}$ and $A = [F_{ijk}]_l -
\Pi_l$:
\begin{align*}
    \big\|[F_{ijk}^{\,\mathrm{gauge}}]_l - \Pi_l\big\|_F
    &= \big\|g_k^{(l)}\, \big([F_{ijk}]_l - \Pi_l\big)\,
       (g_k^{(l)})^{-1}\big\|_F \\
    &= \big\|[F_{ijk}]_l - \Pi_l\big\|_F.
\end{align*}
This holds for all $l \in K$, establishing the full gauge invariance of the
block-specific curvature norms.

\medskip\noindent\textbf{Remark on physical interpretation.} The gauge invariance proved
here ensures that the diagnostic framework of Definition~\ref{def:binding-failure}---which
classifies standpoint failures by examining which $\|F_{ijk}\|_k$ is large---is
well-defined independently of the choice of orthonormal basis within each standpoint
subspace $W_k$. A different gauge choice (a different basis for $W_k$, implemented by an
element of $O(d_k)$) yields the same curvature norms. The diagnosis is a property of the
underlying geometric structure, not an artifact of the representation.
\end{proof}

% ==============================================================================
\subsection{Proof of Proposition: Block Decomposition of Total Curvature}
\label{app:proof-block-decomposition}

\begin{proposition}[Block Decomposition of Total Curvature, Orthogonal Case, restated from
Proposition~\ref{prop:curvature-block-decomposition}]
\label{prop:curvature-block-decomposition-app}
In the orthogonal special case ($\delta = 0$, so $\Pi_k = P_{W_k}$), the total curvature
norm decomposes as
\[
    \|F_{ijk} - I_d\|_F^2 = \sum_{k \in K} \big\|[F_{ijk}]_k - \Pi_k\big\|_F^2
    + \sum_{\substack{k,l \in K \\ k \neq l}} \big\|P_{W_k}\, F_{ijk}\, P_{W_l}\big\|_F^2.
\]
In the general non-orthogonal case ($\delta > 0$), an $O(\delta)$ correction from
cross-layer terms is present (Remark~\ref{rem:block-decomposition}).
\end{proposition}

\begin{proof}
We use the orthogonal decomposition $\mathbb{R}^d = \bigoplus_{k \in K} W_k$ to express
any $d \times d$ matrix as a sum of blocks indexed by pairs $(k, l) \in K \times K$, and
then apply the Pythagorean theorem for the Frobenius inner product.

\medskip\noindent\textbf{Step 1: Block decomposition of a matrix.}
For any matrix $M \in \operatorname{Mat}(d, \mathbb{R})$ and the orthogonal decomposition
$\mathbb{R}^d = \bigoplus_k W_k$ with projectors $\{P_{W_k}\}_{k \in K}$ satisfying
$P_{W_k} P_{W_l} = \delta_{kl}\, P_{W_k}$ and $\sum_k P_{W_k} = I_d$, we can write
\[
    M = I_d\, M\, I_d = \Big(\sum_{k \in K} P_{W_k}\Big)\, M\,
    \Big(\sum_{l \in K} P_{W_l}\Big) = \sum_{k, l \in K} P_{W_k}\, M\, P_{W_l}.
\]
Define the $(k, l)$-block of $M$:
\[
    [M]_{kl} = P_{W_k}\, M\, P_{W_l} \in \operatorname{Hom}(W_l, W_k).
\]
Note that $[M]_{kl}$ maps $W_l$ into $W_k$, and $[M]_{kl} = 0$ when restricted to
$W_l^\perp$.

\medskip\noindent\textbf{Step 2: Orthogonality of blocks in the Frobenius inner product.}
We claim that the blocks $\{[M]_{kl}\}$ are mutually orthogonal in the Frobenius inner
product $\langle A, B \rangle_F = \operatorname{tr}(A^\top B)$. That is, for $(k, l) \neq
(k', l')$:
\[
    \langle [M]_{kl},\, [M]_{k'l'}\rangle_F = 0.
\]
To verify this:
\begin{align*}
    \langle [M]_{kl},\, [M]_{k'l'}\rangle_F
    &= \operatorname{tr}\!\big([M]_{kl}^\top\, [M]_{k'l'}\big) \\
    &= \operatorname{tr}\!\big((P_{W_k}\, M\, P_{W_l})^\top\, P_{W_{k'}}\, M\,
       P_{W_{l'}}\big) \\
    &= \operatorname{tr}\!\big(P_{W_l}\, M^\top\, P_{W_k}\, P_{W_{k'}}\, M\,
       P_{W_{l'}}\big)
       && \text{(since } P_{W_k}^\top = P_{W_k}\text{)} \\
    &= \operatorname{tr}\!\big(P_{W_l}\, M^\top\, \delta_{kk'}\, P_{W_k}\, M\,
       P_{W_{l'}}\big)
       && \text{(since } P_{W_k} P_{W_{k'}} = \delta_{kk'}\, P_{W_k}\text{)} \\
    &= \delta_{kk'}\, \operatorname{tr}\!\big(P_{W_l}\, M^\top\, P_{W_k}\, M\,
       P_{W_{l'}}\big).
\end{align*}
By the cyclic property of trace:
\begin{align*}
    \operatorname{tr}\!\big(P_{W_l}\, M^\top\, P_{W_k}\, M\, P_{W_{l'}}\big)
    &= \operatorname{tr}\!\big(M^\top\, P_{W_k}\, M\, P_{W_{l'}}\, P_{W_l}\big) \\
    &= \delta_{l'l}\, \operatorname{tr}\!\big(M^\top\, P_{W_k}\, M\, P_{W_l}\big),
    && \text{(since } P_{W_{l'}} P_{W_l} = \delta_{l'l}\, P_{W_l}\text{)}
\end{align*}
Therefore:
\[
    \langle [M]_{kl},\, [M]_{k'l'}\rangle_F = \delta_{kk'}\, \delta_{ll'}\,
    \operatorname{tr}\!\big(M^\top\, P_{W_k}\, M\, P_{W_l}\big),
\]
which is zero whenever $(k, l) \neq (k', l')$.

\medskip\noindent\textbf{Step 3: Pythagorean decomposition of the Frobenius norm.}
By the orthogonality of blocks established in Step 2, the Frobenius norm squared of $M$ is
the sum of the Frobenius norms squared of its blocks:
\begin{equation}
\label{eq:pythagorean}
    \|M\|_F^2 = \sum_{k, l \in K} \|[M]_{kl}\|_F^2.
\end{equation}
To verify directly:
\begin{align*}
    \|M\|_F^2
    &= \operatorname{tr}(M^\top M) \\
    &= \operatorname{tr}\!\bigg(\Big(\sum_{k,l} P_{W_l}\, M^\top\, P_{W_k}\Big)\,
       \Big(\sum_{k',l'} P_{W_{k'}}\, M\, P_{W_{l'}}\Big)\bigg) \\
    &= \sum_{k,l,k',l'} \operatorname{tr}\!\big(P_{W_l}\, M^\top\, P_{W_k}\,
       P_{W_{k'}}\, M\, P_{W_{l'}}\big) \\
    &= \sum_{k,l,k',l'} \delta_{kk'}\, \operatorname{tr}\!\big(P_{W_l}\, M^\top\,
       P_{W_k}\, M\, P_{W_{l'}}\big) \\
    &= \sum_{k,l,l'} \delta_{ll'}\, \operatorname{tr}\!\big(M^\top\, P_{W_k}\, M\,
       P_{W_l}\big)
       && \text{(cyclic trace and } P_{W_{l'}} P_{W_l} = \delta_{l'l} P_{W_l}\text{)} \\
    &= \sum_{k,l} \operatorname{tr}\!\big(M^\top\, P_{W_k}\, M\, P_{W_l}\big) \\
    &= \sum_{k,l} \operatorname{tr}\!\big((P_{W_k}\, M\, P_{W_l})^\top\,
       (P_{W_k}\, M\, P_{W_l})\big) \\
    &= \sum_{k,l} \|[M]_{kl}\|_F^2.
\end{align*}

\medskip\noindent\textbf{Step 4: Application to $F_{ijk} - I_d$.}
Set $M = F_{ijk} - I_d$. The blocks of $M$ are:
\[
    [M]_{kl} = P_{W_k}\, (F_{ijk} - I_d)\, P_{W_l}
    = P_{W_k}\, F_{ijk}\, P_{W_l} - P_{W_k}\, I_d\, P_{W_l}.
\]
For the identity term, $P_{W_k}\, I_d\, P_{W_l} = P_{W_k}\, P_{W_l} = \delta_{kl}\,
P_{W_k}$. Therefore:
\[
    [M]_{kl} =
    \begin{cases}
        P_{W_k}\, F_{ijk}\, P_{W_k} - P_{W_k} = [F_{ijk}]_k - I_{d_k}
        & \text{if } k = l, \\[4pt]
        P_{W_k}\, F_{ijk}\, P_{W_l}
        & \text{if } k \neq l,
    \end{cases}
\]
where we identify $[F_{ijk}]_k = P_{W_k}\, F_{ijk}\, P_{W_k}$ as the $k$-th diagonal
block of $F_{ijk}$, and $P_{W_k} = I_{d_k}$ (as an operator on $W_k$). In the orthogonal
case, $\Pi_k = P_{W_k}$, so $[F_{ijk}]_k - \Pi_k = [F_{ijk}]_k - P_{W_k}$.

Applying the Pythagorean decomposition~\eqref{eq:pythagorean}:
\begin{align*}
    \|F_{ijk} - I_d\|_F^2
    &= \sum_{k, l \in K} \|[M]_{kl}\|_F^2 \\
    &= \sum_{k \in K} \|[M]_{kk}\|_F^2 + \sum_{\substack{k, l \in K \\ k \neq l}}
       \|[M]_{kl}\|_F^2 \\
    &= \sum_{k \in K} \big\|[F_{ijk}]_k - I_{d_k}\big\|_F^2
    + \sum_{\substack{k, l \in K \\ k \neq l}} \big\|P_{W_k}\, F_{ijk}\, P_{W_l}\big\|_F^2.
\end{align*}
This is the stated identity.

\medskip\noindent\textbf{Interpretation.} The first sum captures the \emph{diagonal}
(block-specific) curvature: the failure of each standpoint layer's transport to be
trivial. The second sum captures the \emph{off-diagonal} (cross-layer) curvature: the
leakage of transport between distinct standpoint layers. Under
the orthogonal special case ($\delta = 0$) of Assumption~\ref{ass:independence} and the
block-diagonal structure of $U_{ij}$ (Proposition~\ref{prop:block-structure}), the
cross-layer terms vanish at the single-layer level. In the general non-orthogonal case
($\delta > 0$), the Pythagorean decomposition acquires an $O(\delta)$ correction from
cross-layer terms, but the block-specific norms remain the primary diagnostic quantity.
As shown in Proposition~\ref{prop:cross-term}, multi-layer composition can introduce
non-zero off-diagonal curvature, reflecting the dynamic non-commutativity between
standpoint layers (Remark~\ref{rem:non-abelian-source}).
\end{proof}

% ==============================================================================
\subsection{Proof of Theorem: Curvature as Binding Failure --- Formal Equivalence}
\label{app:proof-binding-failure}

We now prove the central theorem connecting the geometric curvature of the fiber bundle to
the binding failure concept of the commitment architecture~\cite{self-jurisdiction}. This
theorem establishes that the two frameworks---differential geometry and commitment
logic---are formally equivalent descriptions of the same phenomenon.

\begin{theorem}[Curvature as Binding Failure: Formal Equivalence, restated from
Theorem~\ref{thm:curvature-binding}]
\label{thm:curvature-binding-app}
Let $\Psi_{x_i}$ be the standpoint state at event $x_i$, and let $U_{ij} \cdot U_{jk}$ be
the composite transport through event $x_j$ to event $x_k$. Then for each layer $k \in K$:
\[
    \|F_{ijk}\|_k > 0 \quad\Longleftrightarrow\quad \text{the commitment encoded in
    layer $k$ was not preserved through event $x_j$.}
\]
\end{theorem}

\begin{proof}
The proof has three parts. First, we establish the formal correspondence between the
commitment architecture and the geometric framework. Second, we prove the forward direction
(positive curvature implies binding failure). Third, we prove the reverse direction (binding
failure implies positive curvature).

\medskip\noindent\textbf{Part 1: Formal correspondence.}
We recall the correspondence between the commitment architecture of~\cite{self-jurisdiction}
and the geometric framework (Table~\ref{tab:correspondence}):
\begin{enumerate}
    \item The self-jurisdiction state $M_t$ corresponds to the standpoint state $\Psi_{x_t}
    = (\psi_k(x_t))_{k \in K}$.
    \item The binding operation $\mathrm{Bind}(c_t, a_t, S_t, \ell_t)$ corresponds to
    connection transport $U_{ij}$.
    \item The commitment in layer $k$ at event $x_i$ is encoded in $\psi_k(x_i) =
    \Pi_k\, h_{x_i}$, the oblique projection of the residual stream onto $W_k$.
    \item The commitment is \emph{preserved} through event $x_j$ if and only if the
    standpoint component $\psi_k$ at event $x_k$ equals its expected value
    $\psi_k^{\exp}(x_k)$, i.e., the value it would have had under flat (standpoint-preserving)
    transport.
\end{enumerate}

\medskip\noindent\textbf{Part 2: Forward direction ($\Rightarrow$).}
Suppose $\|F_{ijk}\|_k > 0$. By definition of the block-specific curvature norm
(Definition~\ref{def:binding-failure}):
\[
    \|F_{ijk}\|_k = \big\|[F_{ijk}]_k - \Pi_k\big\|_F > 0,
\]
which implies $[F_{ijk}]_k \neq \Pi_k$.

Since $U_{ij}$ and $U_{jk}$ are block-diagonal (Proposition~\ref{prop:block-structure}),
their product $U_{ij} \cdot U_{jk}$ is also block-diagonal, with $k$-th block
\[
    [U_{ij} \cdot U_{jk}]_k = \bar{\alpha}_{ji}^{(k)}\, \bar{\alpha}_{kj}^{(k)}\,
    I_{d_k}.
\]
Similarly, the expected transport has $k$-th block $[U_{ik}^{\exp}]_k =
\bar{\alpha}^{(k)}_{\exp}\, I_{d_k}$. The $k$-th block of the curvature is:
\[
    [F_{ijk}]_k = [U_{ij} \cdot U_{jk}]_k \cdot ([U_{ik}^{\exp}]_k)^{-1}
    = \frac{\bar{\alpha}_{ji}^{(k)}\, \bar{\alpha}_{kj}^{(k)}}
           {\bar{\alpha}^{(k)}_{\exp}}\, I_{d_k}.
\]
The condition $[F_{ijk}]_k \neq \Pi_k$ is equivalent to
\[
    \bar{\alpha}_{ji}^{(k)}\, \bar{\alpha}_{kj}^{(k)} \neq \bar{\alpha}^{(k)}_{\exp}.
\]
This means the composite attention weight for layer $k$ through event $x_j$ differs from
the expected value. In the commitment architecture language, the standpoint component
$\psi_k$ was transported to a different value than expected:
\[
    \psi_k(x_k) = [U_{ij} \cdot U_{jk}]_k\, \psi_k(x_i)
    \neq [U_{ik}^{\exp}]_k\, \psi_k(x_i) = \psi_k^{\exp}(x_k).
\]
The event $x_j$ (a pressure turn, such as a challenge or suggestion to revise) induced a
change in $\psi_k$ that was not present in the expected (baseline) transport. By the
correspondence $U_{ij} \leftrightarrow \mathrm{Bind}$ (Table~\ref{tab:correspondence}),
this discrepancy constitutes a binding failure: the commitment encoded in layer $k$ at
event $x_i$ was not preserved through the intervening event $x_j$.

\medskip\noindent\textbf{Part 3: Reverse direction ($\Leftarrow$).}
Suppose the commitment encoded in layer $k$ was not preserved through event $x_j$. By the
formal definition of binding failure in the commitment architecture, this means that the
self-jurisdiction state at event $x_k$ does not cohere with the committed state from
$x_i$, when mediated by the pressure event $x_j$.

Translating to the geometric framework via the correspondence of Part~1: the standpoint
component $\psi_k(x_k)$ differs from its expected value $\psi_k^{\exp}(x_k)$. That is:
\[
    \psi_k(x_k) \neq \psi_k^{\exp}(x_k).
\]
Since $\psi_k(x_k) = [U_{ij} \cdot U_{jk}]_k\, \psi_k(x_i)$ and $\psi_k^{\exp}(x_k) =
[U_{ik}^{\exp}]_k\, \psi_k(x_i)$, this implies
\[
    [U_{ij} \cdot U_{jk}]_k\, \psi_k(x_i) \neq [U_{ik}^{\exp}]_k\, \psi_k(x_i).
\]
Since this must hold for some standpoint state $\psi_k(x_i) \neq 0$ (the commitment is
non-trivial), we conclude
\[
    [U_{ij} \cdot U_{jk}]_k \neq [U_{ik}^{\exp}]_k,
\]
and therefore
\[
    [F_{ijk}]_k = [U_{ij} \cdot U_{jk}]_k \cdot ([U_{ik}^{\exp}]_k)^{-1} \neq \Pi_k.
\]
It follows that
\[
    \|F_{ijk}\|_k = \big\|[F_{ijk}]_k - \Pi_k\big\|_F > 0.
\]

\medskip\noindent\textbf{Conclusion.}
Combining Parts~2 and~3:
\[
    \|F_{ijk}\|_k > 0
    \quad\Longleftrightarrow\quad
    \text{the commitment encoded in layer $k$ was not preserved through event $x_j$.}
\]
This establishes the formal equivalence between the geometric curvature
$\|F_{ijk}\|_k$ and the binding failure of the commitment architecture. The curvature
norm provides a continuous, computable measure of the discrete binary concept of binding
failure: $\|F_{ijk}\|_k = 0$ corresponds to perfect binding (commitment preserved), and
$\|F_{ijk}\|_k > 0$ quantifies the severity of the failure.

\medskip\noindent\textbf{Remark on the low-curvature attractor.}
Theorem~\ref{thm:curvature-binding-app}, combined with the block decomposition of total
curvature (Proposition~\ref{prop:curvature-block-decomposition-app}), shows that the
low-curvature attractor $\mathcal{A}_{\mathrm{self}} = \{\Psi : \|F_\Psi\| < \epsilon\}$
is precisely the set of standpoint configurations in which no binding failure is severe:
for every consistency obligation $\langle x_i, x_j, x_k \rangle$ and every layer $k$, the
commitment encoded in $W_k$ was approximately preserved through the intervening event $x_j$.
Membership in $\mathcal{A}_{\mathrm{self}}$ is gauge-invariant
(Proposition~\ref{prop:gauge-invariance-app}), well-defined
(Theorem~\ref{thm:path-dependence-app}), and computable from attention patterns
(Definition~\ref{def:curvature}). This completes the formal bridge from the geometric
framework to the commitment architecture.
\end{proof}

% ==============================================================================
\section*{Appendix B: Value Matrix Orthonormality Measurement}
\label{app:orthonorm}

The derivation of the block-diagonal transport structure
(Proposition~\ref{prop:block-structure}) requires $V_h^\top V_h = \mathrm{Id}_{d_v}$,
i.e., the columns of each value matrix $V_h$ are orthonormal. This is not enforced by
transformer training. We assess this assumption empirically by computing, for each
attention head $h$, the deviation
\[
    \epsilon_h = \|V_h^\top V_h - \mathrm{Id}_{d_v}\|_F.
\]

\paragraph{GPT-2 (144 heads).}
For GPT-2's 144 attention heads (12 layers $\times$ 12 heads, $d_v = 64$):
\begin{itemize}
    \item Mean $\bar{\epsilon}$: [to be filled after measurement]
    \item Median $\tilde{\epsilon}$: [to be filled]
    \item Maximum $\max_h \epsilon_h$: [to be filled]
\end{itemize}

\paragraph{Llama-2-7b-Chat (1024 heads).}
For Llama-2-7b-Chat's 1024 attention heads (32 layers $\times$ 32 heads, $d_v = 128$):
\begin{itemize}
    \item Mean $\bar{\epsilon}$: [to be filled after measurement]
    \item Median $\tilde{\epsilon}$: [to be filled]
    \item Maximum $\max_h \epsilon_h$: [to be filled]
\end{itemize}

\paragraph{Interpretation.}
When $\bar{\epsilon} < 0.1$, the orthonormality assumption is approximately satisfied and
the derived block-diagonal structure holds to good approximation. When $\bar{\epsilon} >
0.3$, significant off-diagonal cross-block contributions arise beyond those captured by
the non-orthogonality parameter $\delta$. In either case, all holonomy computations in
the main text use the empirically corrected projectors: rather than assuming $V_h V_h^\top
= P_{\mathrm{range}(V_h)}$, we compute the actual projector via SVD of $V_h$ and use the
resulting orthogonal basis for the subspace $W_k$.

\paragraph{Computational note.}
This measurement requires only the value projection matrices $V_h$, which are available
from the pretrained model weights. No forward passes or conversation data are needed. The
computation is $O(H \cdot d_v^3)$ (one SVD per head) and runs in seconds on CPU.
